\chapter*{符号说明}

本书中我们使用了许多数学符号。我们已尽量
保持这些符号的一致性。

作为第一个例子，对于自然数 \(n\)，集合 \(\{1, 2,
\dots, n\}\) 记为 \([n]\)。

\section*{线性代数符号}
\begin{itemize}
    \item 标量（确定性或随机的）使用非粗体表示，
        例如 \(x\)、\(X\)。
    \item 向量（确定性或随机的）使用小写
        粗体表示，例如 \(\vx\)、\(\vpi\)。向量的元素记为
        \(x_{i}\) 或 \(\pi_{i}\)，或者记为 \((\vx)_{i}\)。
        通用项（可以是标量、向量、矩阵或
        高阶张量）也使用小写粗体表示。
    \item 矩阵（确定性或随机的）使用大写粗体表示，
        例如 \(\vX, \vPi\)。矩阵的元素记为
        \(X_{ij}\) 或 \(\Pi_{ij}\)，或者记为 \((\vX)_{ij}\)。
        矩阵的列使用小写粗体向量表示，例如，
        除非另有定义，\(\vx_{i}\) 表示 \(\vX\) 的第 \(i\th\) 列。
    \item 矩阵的转置记为 \(\top\)，例如
        \(\vA^{\top}\)。伴随矩阵记为 \(\adj\)，例如 \(\vA\adj\)。
        伪逆记为 \(\dagger\)，例如 \(\vA^{\dagger}\)。
    \item 矩阵的秩记为 \(\rank(\vA)\)，迹记为
        \(\tr(\vA)\)，行列式记为 \(\det(\vA)\)，
        对数行列式记为 \(\logdet(\vA)\)。
    \item 矩阵的逐元素乘法记为 \(\vA \hada
        \vB\)。逐元素平方记为
        \(\vA^{\hada 2}\)，依此类推。类似地，矩阵的克罗内克积
        记为 \(\vA \kron \vB\)。迭代克罗内克积
        记为 \(\vA^{\kron 2}\)，依此类推。
    \item 对于函数 \(f \colon \R \to \R\)，其对矩阵 \(\vX\) 的逐元素
        应用记为 \(f[\vX]\)，即使用
        方括号。对于函数 \(f \colon \R^{d} \to
        \R^{k}\) 和矩阵 \(\vX \in \R^{d \times n}\)，我们可以
        将 \(f\) 广播（broadcast）到每一列而无需特殊
        符号，即 \(f(\vX) = [f(\vx_{1}), \dots, f(\vx_{n})]
        \in \R^{k \times n}\)。
    \item （欧几里得）单位球面记为 \(\Sphere^{n - 1} \subseteq
        \R^{n}\)。（欧几里得）单位球体记为 \(\Ball^{n} \subseteq \R^{n}\)。
    \item 正交矩阵的集合记为 \(\O(m, n) \subseteq
        \R^{m \times n}\) 或 \(\O(n) = \O(n, n)\)。
    \item 对称矩阵记为 \(\Sym(n)\)，对称半正定（PSD）矩阵
        记为 \(\PSD(n)\)，对称正定（PD）矩阵记为 \(\PD(n)\)。
    \item 根据谱定理，对称矩阵 \(\vS \in \Sym(n)\) 的特征值
        均为实数；它们记为
        \(\lambda_{i}(\vS)\)，并按 \(\lambda_{1}(\vS)
        \geq \cdots \geq \lambda_{n}(\vS)\) 的顺序排列。
    \item 秩为 \(\rank(\vA) = r\) 的矩阵 \(\vA \in \R^{m \times n}\) 的
        奇异值记为 \(\sigma_{i}(\vA)\)，并
        按 \(\sigma_{1}(\vA) \geq \cdots \geq
        \sigma_{r}(\vA) > 0\) 的顺序排列。按照惯例，我们设 \(\sigma_{r +
        1}(\vA) = \sigma_{r + 2}(\vA) = \cdots = 0\)。
    \item 到集合 \(\cK\) 上的投影记为 \(\proj_{\cK}\)。
    \item 向量的 \(\ell^{p}\) 范数记为
        \(\norm{\vx}_{p}\)，其中 \(p \in [0, \infty]\)。
    \item 矩阵的欧几里得算子范数记为 \(\norm{\vA} =
        \sigma_{1}(\vA)\)。
    \item 矩阵的弗罗贝尼乌斯（Frobenius）范数记为 \(\norm{\vA}_{F} =
        \tr(\vA^{\top}\vA)\)。
    \item 向量的欧几里得内积记为 \(\ip{\vx}{\vy} =
        \vy^{\top} \vx\)。
    \item 矩阵的弗罗贝尼乌斯内积记为 \(\ip{\vX}{\vY}
        = \tr(\vY^{\top}\vX)\)。
    \item 全一向量/矩阵记为 \(\vone\)（其形状
        应可从上下文中明显看出）。类似地，全零
        向量/矩阵记为 \(\vzero\)。单位矩阵
        记为 \(\vI\)。
\end{itemize}

\section*{概率符号}

\begin{itemize}
    \item 基础概率测度记为 \(\Pr\)，即
        事件 \(A\) 发生的概率为 \(\Pr[A]\)。我们可以
        使用下标来指定随机变量的分布，
        即 \(\Pr_{\vx \sim \mu}[\vx \in S]\) 描述了
        与服从分布（测度）$\mu$ 的随机变量 $\vx$ 相关的概率。
        如果两个随机变量 $\vx$ 和 $\vx'$ 具有相同的
        分布，我们记为 $\vx \equid \vx'$。
    \item 期望算子记为 \(\Ex\)，下标的用法
        与上述相同。
    \item 协方差算子记为 \(\Cov\)，下标的用法与上述相同。
    \item 相关系数算子记为 \(\Corr\)，下标的用法与上述相同。
    \item 在假设存在密度的某些概率建模环境中
        （特别是在
        \Cref{ch:denoising,ch:conditional-inference} 中），我们也使用符号
        $p_{\vx}$ 来表示随机变量 $\vx$ 的密度。这可以扩展到
        条件密度，例如 $p_{\vx \mid \vy}$ 表示在给定 $\vy$ 的条件下 $\vx$
        的密度。
    \item 我们通常使用希腊字母来表示
        随机变量的实现（以区别于
        随机变量本身）。例如，$p_{\vx}(\vxi)$、$p_{\vx \mid
        \vy}(\vxi \mid \vnu)$、$\Ex[\vx \mid \vy=\vnu]$ 等。
    \item 有限集合 %
        %\footnote{This
        %    should be a measurable space to avoid measure-theoretic and
        %    set-theoretic issues. If you have not studied measure theory,
        %you may relatively safely ignore this.} 
        \(\cX\) 上的概率分布集合记为
        \(\Delta(\cX)\)。（对于 \(\cX = [n]\)，这是可以
        嵌入到 \(\R^{n}\) 中的概率单纯形）。
    \item 均值为 \(\vmu\) 且协方差为
        \(\vSigma\) 的高斯分布记为 \(\dNorm(\vmu, \vSigma)\)。
    \item 紧集 \(\cX\) 上的均匀分布
        记为 \(\dUnif(\cX)\)。
\end{itemize}

\section*{机器学习符号}

\begin{itemize}
    \item 编码器和去噪器通常记为 \(f\)，并且
        通常具有参数 \(\theta \in \Theta\)。我们通常
        将特征写为 \(\vz = f_{\theta}(\vx)\)。
    \item 解码器通常记为 \(g\)，并且通常具有
        参数 \(\eta\)。我们通常将对应于 \(\vx\) 的自编码
        写为 \(\hat{\vx} = g_{\eta}(\vz)\)。
    \item 当 \(f\) 和 \(g\) 实现为神经网络时，
        不失一般性，它们将具有相同的层数；
        我们将其写为 \(f = f^{L} \circ f^{L - 1}
        \circ \cdots \circ f^{2} \circ f^{1} \circ f^{\pre}\) 和 \(g
            = g^{\post} \circ g^{L} \circ g^{ L - 1} \circ \cdots \circ
        g^{2} \circ g^{1}\)。
    \item 我们将第 \(\ell\) 层输入处的特征写为
        \(\vz^{\ell}\)，使得 \(\vz^{\ell + 1} =
        f^{\ell}(\vz^{\ell})\)，其中 \(\vz^{1} = f^{\pre}(\vx)\) 且
        \(\vz = \vz^{L + 1}\)。类似地，自编码特征
        为 \(\hat{\vx}^{\ell}\)，使得 \(\hat{\vx}^{\ell + 1} =
        g^{\ell}(\hat{\vx}^{\ell})\)，其中 \(\hat{\vx}^{1} = \vz\) 且
        \(\hat{\vx} = g^{\post}(\hat{\vx}^{L + 1})\)。
    \item 对于去噪、优化以及具有
        连续时间索引的过程，时间几乎总是作为下标，
        例如 \(\vx_{t}\)。离散序列的索引可以是
        上标或下标。
\end{itemize}